\documentclass[conference]{IEEEtran}
\usepackage{graphicx,booktabs,amsmath,hyperref}
\begin{document}

\title{Interpretable Multi-Signal Claim Verification for Detecting Hallucinations in LLM Answers: A Frozen-Evidence Study}
\author{\IEEEauthorblockN{Your Name}
\IEEEauthorblockA{Department, College\\ City, Country\\ email@example.com}}
\maketitle

\begin{abstract}
LLMs produce fluent but unsupported statements. We present a lightweight, fully local, interpretable
pipeline that decomposes answers into claims and assigns one of four verdicts (Verified, Partially
Supported, Unsupported, Contradicted) using subject-centred retrieval, evidence-quality ranking,
answer-value checking, entity/relationship checks and NLI. Using a frozen-evidence protocol that
isolates verification from retrieval and generation, we report ablations, baselines, confidence
analysis and a failure taxonomy on a hand-built challenge set and a VitaminC-derived benchmark.
\textbf{[Insert 1--2 headline numbers with CIs.]}
\end{abstract}

\begin{IEEEkeywords}
hallucination detection, fact verification, natural language inference, evidence retrieval
\end{IEEEkeywords}

\section{Introduction}
Motivation; the wrong-entity retrieval failure (e.g., a birthplace claim ``verified'' by the article
about the claimed city); contributions: (1) subject-centred retrieval with a subject-presence guard,
(2) four-way verdicts with explicit insufficient-evidence semantics and a false-VERIFIED safety metric,
(3) a frozen-evidence evaluation protocol, (4) ablations, confidence analysis and a failure taxonomy.

\section{Related Work}
FEVER \cite{fever}, VitaminC \cite{vitaminc}, NLI-based verification \cite{mnli,bart},
sentence embeddings \cite{sbert}, hallucination detection \cite{selfcheck}.

\section{Method}
Pipeline figure; question analysis; subject-centred queries; quality score
$q = 0.30\,s + 0.30\,u + 0.20\,e + 0.10\,r + 0.10\,w$; verdict rules 0--7 (table);
heuristic confidence $0.35\,s+0.25\,a+0.20\,c+0.20\,e$ (uncalibrated by design).

\section{Experimental Setup}
Datasets: challenge set (40 hand-built claims), VitaminC subset (dev/test, grouped split by claim).
Protocol: frozen evidence; thresholds tuned on dev only; test evaluated once.
Metrics: accuracy, macro-F1 over populated classes, false-VERIFIED rate, bootstrap 95\% CIs (1000 resamples).
Baselines: similarity+agreement, NLI only (bart-large-mnli), Gemini judge, logistic-regression combiner.

\section{Results}
\input{../benchmark/results/vc_test_ablation.tex}
\textbf{Main result / ablation} -- discuss each added component's contribution.

\begin{table}[t]
\centering
\caption{Ablation on challenge (frozen evidence).}
\label{tab:ablation_challenge}
\small
\begin{tabular}{lccc}
\hline
Configuration & Acc [95\\% CI] & Macro-F1 & False-VER \\
\hline
Sim+Agreement & 0.375 [0.23,0.53] & 0.506 & 0.167 \\
NLI only & 0.700 [0.57,0.85] & 0.671 & 0.033 \\
Sim+NLI & 0.675 [0.53,0.80] & 0.714 & 0.067 \\
Sim+NLI+Entity & 0.675 [0.53,0.80] & 0.714 & 0.067 \\
Sim+NLI+Entity+Rel & 0.575 [0.42,0.72] & 0.665 & 0.067 \\
Full (+AnswerValue) & 0.575 [0.42,0.72] & 0.665 & 0.067 \\
Engine (determine\_verdict) & 0.575 [0.42,0.72] & 0.665 & 0.067 \\
\hline
\end{tabular}
\end{table}
\textbf{Challenge set} -- confusion matrix and safety metric.

\textbf{Confidence analysis} -- AUROC, ECE, reliability (Fig.~\ref{fig:rel}).
\begin{figure}[t]\centering
\includegraphics[width=0.8\columnwidth]{../benchmark/results/conf_reliability.png}
\caption{Reliability diagram of the heuristic confidence.}\label{fig:rel}
\end{figure}

\input{../benchmark/results/challenge_errors.tex}
\textbf{Error analysis} -- 3--4 concrete examples per category.

\section{Limitations}
Small scale; single-evidence VitaminC; heuristic thresholds; spaCy NER dependence; English/Wikipedia only;
challenge-set entities overlap with development entities.

\section{Conclusion}
Summary and future work (learned verifier, larger multi-evidence benchmarks, live-retrieval evaluation).

\begin{thebibliography}{9}
\bibitem{fever} J. Thorne et al., ``FEVER: a large-scale dataset for fact extraction and verification,'' in \emph{NAACL}, 2018.
\bibitem{vitaminc} T. Schuster, A. Fisch, R. Barzilay, ``Get your Vitamin C! Robust fact verification with contrastive evidence,'' in \emph{NAACL}, 2021.
\bibitem{mnli} A. Williams, N. Nangia, S. Bowman, ``A broad-coverage challenge corpus for sentence understanding through inference,'' in \emph{NAACL}, 2018.
\bibitem{bart} M. Lewis et al., ``BART: denoising sequence-to-sequence pre-training,'' in \emph{ACL}, 2020.
\bibitem{sbert} N. Reimers, I. Gurevych, ``Sentence-BERT: sentence embeddings using Siamese BERT-networks,'' in \emph{EMNLP-IJCNLP}, 2019.
\bibitem{selfcheck} P. Manakul, A. Liusie, M. Gales, ``SelfCheckGPT: zero-resource black-box hallucination detection,'' in \emph{EMNLP}, 2023.
\end{thebibliography}
\end{document}
